\documentclass[10pt,a4paper]{amsart}
\usepackage{amsmath,amssymb,amsthm,graphicx}
\usepackage{verbatim}
\newenvironment{reference}{\comment}{\endcomment}
%\newenvironment{reference}{}{}
\newenvironment{slogan}{\comment}{\endcomment}
\newenvironment{history}{\comment}{\endcomment}
% For commutative diagrams we use Xy-pic
\usepackage[all]{xy}

% We use 2cell for 2-commutative diagrams.
\xyoption{2cell}
% Theorem environments.
%
\theoremstyle{plain}
\newtheorem{theorem}[subsection]{Theorem}
\newtheorem{proposition}[subsection]{Proposition}
\newtheorem{lemma}[subsection]{Lemma}

\theoremstyle{definition}
\newtheorem{definition}[subsection]{Definition}
\newtheorem{example}[subsection]{Example}
\newtheorem{exercise}[subsection]{Exercise}
\newtheorem{situation}[subsection]{Situation}

\theoremstyle{remark}
\newtheorem{remark}[subsection]{Remark}
\newtheorem{remarks}[subsection]{Remarks}

\numberwithin{equation}{subsection}

% Macros
%
\def\lim{\mathop{\mathrm{lim}}\nolimits}
\def\colim{\mathop{\mathrm{colim}}\nolimits}
\def\Spec{\mathop{\mathrm{Spec}}}
\def\Hom{\mathop{\mathrm{Hom}}\nolimits}
\def\Ext{\mathop{\mathrm{Ext}}\nolimits}
\def\SheafHom{\mathop{\mathcal{H}\!\mathit{om}}\nolimits}
\def\SheafExt{\mathop{\mathcal{E}\!\mathit{xt}}\nolimits}
\def\Sch{\mathit{Sch}}
\def\Mor{\mathop{\mathrm{Mor}}\nolimits}
\def\Ob{\mathop{\mathrm{Ob}}\nolimits}
\def\Sh{\mathop{\mathit{Sh}}\nolimits}
\def\NL{\mathop{N\!L}\nolimits}
\def\CH{\mathop{\mathrm{CH}}\nolimits}
\def\proetale{{pro\text{-}\acute{e}tale}}
\def\etale{{\acute{e}tale}}
\def\QCoh{\mathit{QCoh}}
\def\Ker{\mathop{\mathrm{Ker}}}
\def\Im{\mathop{\mathrm{Im}}}
\def\Coker{\mathop{\mathrm{Coker}}}
\def\Coim{\mathop{\mathrm{Coim}}}

% Boxtimes
%
\DeclareMathSymbol{\boxtimes}{\mathbin}{AMSa}{"02}

%
% Macros for moduli stacks/spaces
%
\def\QCohstack{\mathcal{QC}\!\mathit{oh}}
\def\Cohstack{\mathcal{C}\!\mathit{oh}}
\def\Spacesstack{\mathcal{S}\!\mathit{paces}}
\def\Quotfunctor{\mathrm{Quot}}
\def\Hilbfunctor{\mathrm{Hilb}}
\def\Curvesstack{\mathcal{C}\!\mathit{urves}}
\def\Polarizedstack{\mathcal{P}\!\mathit{olarized}}
\def\Complexesstack{\mathcal{C}\!\mathit{omplexes}}
% \Pic is the operator that assigns to X its picard group, usage \Pic(X)
% \Picardstack_{X/B} denotes the Picard stack of X over B
% \Picardfunctor_{X/B} denotes the Picard functor of X over B
\def\Pic{\mathop{\mathrm{Pic}}\nolimits}
\def\Picardstack{\mathcal{P}\!\mathit{ic}}
\def\Picardfunctor{\mathrm{Pic}}
\def\Deformationcategory{\mathcal{D}\!\mathit{ef}}
\usepackage{cleveref}
\crefname{originalreference}{original reference}{original references}
\Crefname{originalreference}{Original reference}{Original references}
\makeatletter
\newlabel{algebra-lemma-CM-ass-minimal-support}{{algebra-lemma-CM-ass-minimal-support}{source}{}{sourcepointer-1}{}}
\newlabel{algebra-lemma-CM-ass-minimal-support@cref}{{[originalreference][0][]algebra-lemma-CM-ass-minimal-support}{[1][0][]source}{}{}{}}
\makeatother
\crefname{originalreference}{original reference}{original references}
\Crefname{originalreference}{Original reference}{Original references}
\makeatletter
\newlabel{algebra-lemma-ideal-nonzerodivisor}{{algebra-lemma-ideal-nonzerodivisor}{source}{}{sourcepointer-2}{}}
\newlabel{algebra-lemma-ideal-nonzerodivisor@cref}{{[originalreference][0][]algebra-lemma-ideal-nonzerodivisor}{[1][0][]source}{}{}{}}
\makeatother
\crefname{originalreference}{original reference}{original references}
\Crefname{originalreference}{Original reference}{Original references}
\makeatletter
\newlabel{algebra-lemma-one-equation-module}{{algebra-lemma-one-equation-module}{source}{}{sourcepointer-3}{}}
\newlabel{algebra-lemma-one-equation-module@cref}{{[originalreference][0][]algebra-lemma-one-equation-module}{[1][0][]source}{}{}{}}
\makeatother
\crefname{originallemma}{original lemma}{original lemmas}
\Crefname{originallemma}{Original lemma}{Original lemmas}
\makeatletter
\newlabel{lemma-depth-drops-by-one}{{lemma-depth-drops-by-one}{source}{}{sourcepointer-4}{}}
\newlabel{lemma-depth-drops-by-one@cref}{{[originallemma][0][]lemma-depth-drops-by-one}{[1][0][]source}{}{}{}}
\makeatother
\title{The depth of a Cohen-Macaulay module against an ideal as a difference of dimensions}
\author{The Stacks Project Authors}
\date{stacks-project:dualizing.tex:local-snapshot-20261012; GNU FDL 1.2 or later}
\begin{document}
\maketitle
\noindent\textit{Editorial note.} All delivered mathematical text is the author's own, copied from the pinned Stacks Project chapter source sources/stacks-dualizing/pinned-original-source.tex; no statement or proof step was rewritten, reordered or removed.\par\medskip
\noindent\textit{Editorial note.} Preamble provenance: the Stacks Project keeps one single shared preamble for every chapter; this card takes it directly from the staged sources/stacks-dualizing/preamble.tex as three verbatim spans (lines 8-12, 14-18 and 154-217), so no derived file was written and nothing was added or redefined. The wrapper class and the base packages are supplied by the assembler.\par\medskip
\noindent\textit{Editorial note.} Bibliography provenance: neither the counted statement nor its proof contains a citation (this lane rescans the whole span before writing the card), so the card carries no bibliography and no entry was reconstructed.\par\medskip
\noindent\textit{Editorial note.} Selection provenance: this card was opened under the file-based control contract (CONTROL.md), with the mechanical pre-card gate run on its statement and proof interval before the card was written.\par\medskip
\noindent\textit{Allowed standard prerequisites (editorial).}\begin{itemize}
\item The description of the associated primes of a module when the depth against an ideal is zero. Invoked inside the counted proof (source lines 2073-2075). Places the ideal inside an associated prime in the zero depth case.
\item The equality of the dimension of the support with the dimension of the quotient by a minimal associated prime for a Cohen-Macaulay module. Invoked inside the counted proof (source lines 2077-2079). Turns the inequality in the zero depth case into an equality.
\item The drop of the depth by exactly one after dividing by a nonzerodivisor in the ideal. Invoked inside the counted proof (source lines 2086-2087). Provides the induction step on the depth.
\item The drop of the dimension of the support by exactly one after dividing by a nonzerodivisor. Invoked inside the counted proof (source lines 2088-2089). Keeps the two sides of the identity in step during the induction.
\end{itemize}\medskip
\ifdefined\hypertarget\hypertarget{sourcepointer-1}{}\fi
\noindent\textit{Source reference.} Original reference algebra-lemma-CM-ass-minimal-support: sources/stacks-dualizing/pinned-original-source.tex, source lines 2079--2079. Navigation pointer only: the counted text refers to the Stacks item algebra-lemma-CM-ass-minimal-support, which belongs to another chapter of the project and is not part of this card. The pointer carries the target's own label as its number because this chapter's staged snapshot does not contain the other chapter's tag table; its content is not reproduced here.\par\medskip
\ifdefined\hypertarget\hypertarget{sourcepointer-2}{}\fi
\noindent\textit{Source reference.} Original reference algebra-lemma-ideal-nonzerodivisor: sources/stacks-dualizing/pinned-original-source.tex, source lines 2075--2075. Navigation pointer only: the counted text refers to the Stacks item algebra-lemma-ideal-nonzerodivisor, which belongs to another chapter of the project and is not part of this card. The pointer carries the target's own label as its number because this chapter's staged snapshot does not contain the other chapter's tag table; its content is not reproduced here.\par\medskip
\ifdefined\hypertarget\hypertarget{sourcepointer-3}{}\fi
\noindent\textit{Source reference.} Original reference algebra-lemma-one-equation-module: sources/stacks-dualizing/pinned-original-source.tex, source lines 2089--2089. Navigation pointer only: the counted text refers to the Stacks item algebra-lemma-one-equation-module, which belongs to another chapter of the project and is not part of this card. The pointer carries the target's own label as its number because this chapter's staged snapshot does not contain the other chapter's tag table; its content is not reproduced here.\par\medskip
\ifdefined\hypertarget\hypertarget{sourcepointer-4}{}\fi
\noindent\textit{Source reference.} Original lemma lemma-depth-drops-by-one: sources/stacks-dualizing/pinned-original-source.tex, source lines 2039--2039. Navigation pointer only: the counted text refers to the Stacks item lemma-depth-drops-by-one, which is an existing labeled result of this same chapter and is not part of this card. The pointer carries the target's own label as its number and its definition line as its span; its content is not reproduced here.\par\medskip
\begin{lemma}
\label{lemma-depth-CM}
Let $R$ be a Noetherian local ring. If $M$ is a finite Cohen-Macaulay
$R$-module and $I \subset R$ a nontrivial ideal. Then
$$
\text{depth}_I(M) = \dim(\text{Supp}(M)) - \dim(\text{Supp}(M/IM)).
$$
\end{lemma}

\begin{proof}
We will prove this by induction on $\text{depth}_I(M)$.

\medskip\noindent
If $\text{depth}_I(M) = 0$, then $I$ is contained in one
of the associated primes $\mathfrak p$ of $M$
(Algebra, Lemma \ref{algebra-lemma-ideal-nonzerodivisor}).
Then $\mathfrak p \in \text{Supp}(M/IM)$, hence
$\dim(\text{Supp}(M/IM)) \geq \dim(R/\mathfrak p) = \dim(\text{Supp}(M))$
where equality holds by
Algebra, Lemma \ref{algebra-lemma-CM-ass-minimal-support}.
Thus the lemma holds in this case.

\medskip\noindent
If $\text{depth}_I(M) > 0$, we pick $x \in I$ which is a
nonzerodivisor on $M$. Note that $(M/xM)/I(M/xM) = M/IM$.
On the other hand we have
$\text{depth}_I(M/xM) = \text{depth}_I(M) - 1$
by Lemma \ref{lemma-depth-drops-by-one}
and $\dim(\text{Supp}(M/xM)) = \dim(\text{Supp}(M)) -  1$
by Algebra, Lemma \ref{algebra-lemma-one-equation-module}.
Thus the result by induction hypothesis.
\end{proof}

\end{document}
